\subsection{Bifurcation Cascade}
For implementation of the map we use the function \verb+coupled_map+ from \verb+common.py+. This function calculates the next generation
 of x and y values from the given x and y values, $r_1$, $r_2$ and $\epsilon$. As we need more than one time step, we created the 
 function \verb+calc_timeline+, which returns the x and y values for a certain $r_2$ value. This function also has the option to choose 
 different numbers of steps and discard a certain number of timesteps from the beginning of the time series. \\

Additionally, we set our constants to the values given in the task: $r_1 = 3.1$ and $\varepsilon=0.3$. For $r_2$ we created a linspace 
from 2.8 to 3.97 with 800 values. The x and y values each iteration starts from are set to $x_0=0.2$ and $y_0=0.6$. 
We decided to use 400 iterations and discard the first 150. \\
 
For each value in the $r_2$ linspace we then calculate the time series and save the x values. We then use \verb+plt.plot+ to plot all these 
values and then move to the next $r_2$ value. Lastly, we added labels and the title and saved the figure. 
